\documentclass[11pt,a4paper]{article}
\usepackage[utf8]{vietnam}
\usepackage{amsmath,amssymb}
\usepackage{tabularx}
\usepackage{tikz}
\usepackage{xcolor}
\usepackage{fancyhdr}
\usepackage[top=1.8cm, bottom=2.0cm, left=1.8cm, right=1.8cm]{geometry}

\setlength{\parskip}{3pt}
\setlength{\parindent}{0pt}

\pagestyle{fancy}
\fancyhf{}
\renewcommand{\headrulewidth}{0.4pt}
\lhead{\footnotesize \textbf{CHÙA HƯƠNG SƠN -- ÔN THI VÀO 10}}
\rhead{\footnotesize \textbf{ĐÁP ÁN \& LỜI GIẢI CHI TIẾT (ĐỀ SỐ 09)}}
\cfoot{\thepage}

\begin{document}

\begin{center}
\textbf{\large HƯỚNG DẪN CHẤM VÀ ĐÁP ÁN ĐỀ THI THỬ SỐ 09}\\
\textbf{MÔN: TOÁN 9 -- HỌC KỲ I -- NĂM HỌC 2026 -- 2027}\\
\textit{(Hướng dẫn chấm gồm 02 trang)}
\end{center}

\vspace{1mm}
\noindent
\rule{\linewidth}{0.8pt}

\vspace{2mm}
\noindent
\textbf{\large A. ĐÁP ÁN TRẮC NGHIỆM (7,0 ĐIỂM)}

\vspace{2mm}
\noindent
\textbf{1. Phần 1: Trắc nghiệm nhiều lựa chọn (3,0 điểm -- Mỗi câu đúng 0,25 điểm)}
\begin{center}
\begin{tabular}{|c|c|c|c|c|c|c|c|c|c|c|c|}
\hline
\textbf{1} & \textbf{2} & \textbf{3} & \textbf{4} & \textbf{5} & \textbf{6} & \textbf{7} & \textbf{8} & \textbf{9} & \textbf{10} & \textbf{11} & \textbf{12} \\
\hline
\textbf{A} & \textbf{A} & \textbf{A} & \textbf{A} & \textbf{A} & \textbf{A} & \textbf{A} & \textbf{A} & \textbf{A} & \textbf{A} & \textbf{A} & \textbf{A} \\
\hline
\end{tabular}
\end{center}

\vspace{2mm}
\noindent
\textbf{2. Phần 2: Trắc nghiệm Đúng -- Sai (2,0 điểm)}
\begin{itemize}
  \item \textbf{Câu 13 (1,0 điểm):} a) Đúng \quad b) Đúng \quad c) Sai \quad d) Đúng.
  \item \textbf{Câu 14 (1,0 điểm):} a) Đúng \quad b) Đúng \quad c) Sai \quad d) Sai.
\end{itemize}
\textit{Quy tắc chấm: Đúng 1 ý được 0,1đ; đúng 2 ý được 0,25đ; đúng 3 ý được 0,5đ; đúng 4 ý được 1,0đ.}

\vspace{2mm}
\noindent
\textbf{3. Phần 3: Trắc nghiệm trả lời ngắn (2,0 điểm -- Mỗi câu đúng 0,5 điểm)}
\begin{itemize}
  \item \textbf{Câu 15:} $x > 14$.
  \item \textbf{Câu 16:} $4$.
  \item \textbf{Câu 17:} $21{,}7\text{ m}$ (hoặc $21{,}7$).
  \item \textbf{Câu 18:} $141{,}3\text{ cm}^2$ (hoặc $141{,}3$).
\end{itemize}

\vspace{3mm}
\noindent
\textbf{\large B. LỜI GIẢI CHI TIẾT VÀ BAREM ĐIỂM TỰ LUẬN (3,0 ĐIỂM)}

\vspace{2mm}
\noindent
\textbf{Câu 19 (1,0 điểm).}
\begin{itemize}
  \item[a)] \textit{(0,5 điểm)} Hệ phương trình: $\begin{cases} 2x + 5y = 11 \quad (1) \\ 3x - 2y = 7 \quad (2) \end{cases}$.\\
  Nhân (1) với 2 và nhân (2) với 5: $\begin{cases} 4x + 10y = 22 \\ 15x - 10y = 35 \end{cases}$ \dotfill \textbf{(0,25đ)}\\
  Cộng hai phương trình vế theo vế: $19x = 57 \iff x = 3$.\\
  Thay $x = 3$ vào (1): $2(3) + 5y = 11 \iff 5y = 5 \iff y = 1$.\\
  Vậy hệ phương trình có nghiệm duy nhất $(x; y) = (3; 1)$. \dotfill \textbf{(0,25đ)}

  \item[b)] \textit{(0,5 điểm)} Với $x > 0, x \ne 1$, ta có:
  \[
  \dfrac{1}{\sqrt{x}-1} + \dfrac{1}{\sqrt{x}+1} = \dfrac{\sqrt{x}+1 + \sqrt{x}-1}{(\sqrt{x}-1)(\sqrt{x}+1)} = \dfrac{2\sqrt{x}}{x-1} \quad \text{\textbf{(0,25đ)}}
  \]
  Do đó:
  \[
  P = \dfrac{2\sqrt{x}}{x-1} : \dfrac{2\sqrt{x}}{x-1} = \dfrac{2\sqrt{x}}{x-1} \cdot \dfrac{x-1}{2\sqrt{x}} = 1 \quad \text{\textbf{(0,25đ)}}
  \]
\end{itemize}

\vspace{2mm}
\noindent
\textbf{Câu 20 (1,5 điểm).}
\begin{itemize}
  \item[a)] \textit{(0,75 điểm)}
  Vì $MA, MB$ là hai tiếp tuyến của $(O)$ cắt nhau tại $M$ nên $MA = MB$.\\
  Lại có $OA = OB = R$, suy ra $OM$ là đường trung trực của đoạn thẳng $AB$.\\
  Do đó $OM \perp AB$ tại $H$ (đpcm). \dotfill \textbf{(0,5đ)}\\
  Xét tam giác $OAM$ vuông tại $A$, có đường cao $AH \perp OM$.\\
  Áp dụng hệ thức lượng: $MA^2 = MH \cdot MO$ (đpcm). \dotfill \textbf{(0,25đ)}

  \item[b)] \textit{(0,75 điểm)}
  Xét hai tam giác $MAC$ và $MDA$:\\
  Có góc $\widehat{M}$ chung, và $\widehat{MAC} = \widehat{MDA}$ (góc tạo bởi tiếp tuyến và dây cung cùng chắn cung $AC$).\\
  Suy ra $\triangle MAC \backsim \triangle MDA$ (g-g).\\
  Do đó $\dfrac{MA}{MD} = \dfrac{MC}{MA} \implies MA^2 = MC \cdot MD$. \dotfill \textbf{(0,25đ)}\\
  Từ câu a ta có $MA^2 = MH \cdot MO$, suy ra $MC \cdot MD = MH \cdot MO$. \dotfill \textbf{(0,25đ)}\\
  Từ $MC \cdot MD = MH \cdot MO \implies \dfrac{MC}{MO} = \dfrac{MH}{MD}$.\\
  Xét hai tam giác $MCH$ và $MOD$ có góc $\widehat{M}$ chung và $\dfrac{MC}{MO} = \dfrac{MH}{MD}$, nên $\triangle MCH \backsim \triangle MOD$ (c-g-c).\\
  Suy ra $\widehat{MHC} = \widehat{MDO}$. Do đó tứ giác $CHOD$ có góc ngoài bằng góc đối trong, suy ra bốn điểm $C, H, O, D$ cùng thuộc một đường tròn (đpcm). \dotfill \textbf{(0,25đ)}
\end{itemize}

\vspace{2mm}
\noindent
\textbf{Câu 21 (0,5 điểm).}
\begin{itemize}
  \item[a)] \textit{(0,25 điểm)}
  Gọi chiều rộng của mảnh vườn là $x$ ($x > 0$, mét), chiều dài là $y$ ($y > 0$, mét).\\
  Diện tích mảnh vườn: $xy = 200 \implies y = \dfrac{200}{x}$.\\
  Chiều dài hàng rào bao quanh ba phía (gồm 2 chiều rộng và 1 chiều dài) là:
  \[
  L = 2x + y = 2x + \dfrac{200}{x}.
  \]
  Áp dụng bất đẳng thức Cô-si cho hai số dương $2x$ và $\dfrac{200}{x}$:
  \[
  L = 2x + \dfrac{200}{x} \ge 2\sqrt{2x \cdot \dfrac{200}{x}} = 2\sqrt{400} = 2 \cdot 20 = 40\text{ (m)}.
  \]
  Dấu ``='' xảy ra khi $2x = \dfrac{200}{x} \iff 2x^2 = 200 \iff x^2 = 100 \iff x = 10\text{ (m)}$ (do $x > 0$).\\
  Khi đó chiều dài là $y = \dfrac{200}{10} = 20\text{ (m)}$.\\
  Vậy để chiều dài hàng rào là ngắn nhất ($40\text{ m}$) thì chiều rộng là $10\text{ m}$ và chiều dài là $20\text{ m}$. \dotfill \textbf{(0,25đ)}

  \item[b)] \textit{(0,25 điểm)}
  Áp dụng bất đẳng thức Cauchy-Schwarz dạng phân thức (BĐT Engel):
  \[
  Q = \dfrac{a^2}{b+c} + \dfrac{b^2}{c+a} + \dfrac{c^2}{a+b} \ge \dfrac{(a+b+c)^2}{(b+c) + (c+a) + (a+b)} = \dfrac{(a+b+c)^2}{2(a+b+c)} = \dfrac{a+b+c}{2}.
  \]
  Mặt khác, với mọi $a, b, c > 0$ ta luôn có:
  \[
  (a+b+c)^2 \ge 3(ab+bc+ca) = 3 \cdot 3 = 9 \implies a+b+c \ge 3.
  \]
  Suy ra:
  \[
  Q \ge \dfrac{3}{2}.
  \]
  Dấu đẳng thức xảy ra khi $\begin{cases} a = b = c \\ ab+bc+ca=3 \end{cases} \iff a = b = c = 1$.\\
  Vậy giá trị nhỏ nhất của $Q$ là $\dfrac{3}{2}$, đạt được khi $a = b = c = 1$. \dotfill \textbf{(0,25đ)}
\end{itemize}

\end{document}
